% Generated deterministically from the audited Markdown source.
% Responsible author: Carptopus.
\documentclass[11pt]{article}
\usepackage[a4paper,margin=27mm]{geometry}
\usepackage[T1]{fontenc}
\usepackage[utf8]{inputenc}
\usepackage{lmodern}
\usepackage{microtype}
\usepackage{amsmath,amssymb,amsthm,mathtools}
\usepackage{xcolor}
\usepackage[colorlinks=true,linkcolor=blue!55!black,citecolor=blue!55!black,urlcolor=blue!65!black]{hyperref}
\usepackage{xurl}
\usepackage{fancyhdr}

\newtheorem{theorem}{Theorem}
\newtheorem{proposition}[theorem]{Proposition}
\newtheorem{lemma}[theorem]{Lemma}
\newtheorem{corollary}[theorem]{Corollary}

\pagestyle{fancy}
\fancyhf{}
\fancyhead[L]{Carptopus}
\fancyhead[R]{Minor-connectivity ceiling gaps}
\fancyfoot[C]{\thepage}
\setlength{\headheight}{14pt}
\setlength{\parindent}{0pt}
\setlength{\parskip}{0.42em}
\setlength{\emergencystretch}{4em}

\title{An unbounded gap between minimum degree\\
and the minor-connectivity ceiling}
\author{Carptopus\\\href{mailto:carptopus@163.com}{carptopus@163.com}}
\date{3 September 2026}

\hypersetup{
  pdftitle={An unbounded gap between minimum degree and the minor-connectivity ceiling},
  pdfauthor={Carptopus},
  pdfsubject={Exact connectivity ceilings for independent blow-ups of the Mader graph},
  pdfkeywords={graph minor, vertex connectivity, minimum degree, Mader graph, graph blow-up, Barát conjecture}
}

\begin{document}
\maketitle

\begin{center}
\fcolorbox{black!35}{black!4}{\parbox{0.88\textwidth}{\centering\small
Public Beta manuscript, version 0.1-beta. The mathematical claims and reproducibility artifacts
have passed project-internal audits; external mathematical review is pending.}}
\end{center}

\begin{abstract}


For a nonempty finite simple graph $G$, let

\nopagebreak[4]
\[
\kappa^*(G)=\max\{\kappa(H):H\text{ is a nonempty minor of }G\},
\]

where $\kappa$ denotes vertex connectivity. We determine this parameter exactly
on the independent blow-ups of the twelve-vertex Mader graph $M_{12}$. For every
integer $t\geq2$,

\nopagebreak[4]
\[
\kappa^*\bigl(M_{12}[\overline K_t]\bigr)
=\left\lfloor\frac{9t}{2}\right\rfloor,
\qquad
\delta\bigl(M_{12}[\overline K_t]\bigr)=5t.
\]

Consequently, the difference between minimum degree and the largest
connectivity of a minor is unbounded, giving infinitely many counterexamples
to a conjecture of Barát. For context, we also prove the elementary identity
$\kappa^*(K_s\vee G)=s+\kappa^*(G)$; applied to $M_{12}$, it already answers
negatively the question whether sufficiently large minimum degree guarantees
$\kappa^*(G)\geq\delta(G)$, but only with constant deficit. The independent
blow-ups make that deficit unbounded. Their lower bound is supplied by an
explicit uniform minor model. For the upper bound, a two-fibre separator
localizes every relevant minor to a torso that is the join of three elementary
graphs. Separating mixed and pure branch sets then reduces the problem to three
sharp full-quotient profiles and an anticomplete-shores formula for
connectivity. The proof is theoretical; accompanying computations only check
the construction.

\end{abstract}

\section{Introduction}


All ambient graphs in this paper are nonempty, finite and simple; auxiliary
graphs in decompositions may be empty. For a nonempty graph $G$, define the
minor-monotone ceiling of vertex connectivity by

\nopagebreak[4]
\[
\kappa^*(G)=\max\{\kappa(H):H\preccurlyeq G,\ |V(H)|\geq1\}. \tag{1}
\]

where $H\preccurlyeq G$ means that $H$ is a minor of $G$. We use the convention
$\kappa(K_n)=n-1$. The parameter in (1) is denoted $\lceil\kappa\rceil$ in the
minor-monotone ceiling literature.

Barát conjectured that every graph of minimum degree at least $k+1$ contains a
$k$-connected minor. Equivalently, the conjecture asserts

\nopagebreak[4]
\[
\kappa^*(G)\geq\delta(G)-1                                  \tag{2}
\]

for every nonempty graph \cite{barat}. Barioli, Barrett, Fallat, Hall, Hogben, Shader, van
den Driessche and van der Holst studied minor-monotone ceilings and recorded a
twelve-vertex graph $M_{12}$ with

\nopagebreak[4]
\[
\delta(M_{12})=5,
\qquad
\kappa^*(M_{12})=4.                                         \tag{3}
\]

Thus $M_{12}$ attains equality in (2), rather than contradicting it \cite{barioli2013}. More
recently, Barioli, Fallat, Gupta and Li used the same example when asking \cite{barioli2026}
whether there is a universal constant $c$ such that

\nopagebreak[4]
\[
\kappa^*(G)\geq\delta(G)                                    \tag{4}
\]

whenever $\delta(G)\geq c$.

There is a simple first amplification of (3). Proposition~\ref{prop:join} below shows that
joining a graph with a clique increases its connectivity ceiling and minimum
degree by the clique order. Thus $K_s\vee M_{12}$ has arbitrarily large minimum
degree but retains deficit one, already answering (4) negatively. Our main
result is stronger: independent blow-ups of $M_{12}$ have an exact connectivity
ceiling and a deficit that grows linearly.

\begin{theorem}
\label{thm:main}

For every integer $t\geq2$, let

\nopagebreak[4]
\[
G_t=M_{12}[\overline K_t].
\]

Then

\nopagebreak[4]
\[
\kappa^*(G_t)=\left\lfloor\frac{9t}{2}\right\rfloor.        \tag{5}
\]

Since $\delta(G_t)=5t$, Theorem~\ref{thm:main} gives

\nopagebreak[4]
\[
\delta(G_t)-\kappa^*(G_t)=\left\lceil\frac t2\right\rceil. \tag{6}
\]

In particular, (2) fails for every $t\geq3$, and the difference in (6) is
unbounded. This strengthens the fixed-deficit obstruction supplied by clique
joins.

The proof has two complementary parts. The lower bound in (5) comes from a
balanced family of explicit branch sets in the original graph. For the upper
bound, the natural $2t$-vertex separator of $G_t$ localizes any more than
$2t$-connected minor to a $7t$-vertex torso. That torso has a three-factor join
decomposition. Every branch set meeting two factors becomes universal in the
quotient, while branch sets contained in one factor are controlled by simple
full-quotient profiles. The resulting capacity inequalities give the matching upper
bound.

\end{theorem}

\section{A baseline amplification by clique joins}


The ordinary connectivity formula for graph joins is classical. Barrett,
Fallat, Hall and Hogben used it in their study of the connectivity ceiling \cite{barrett2013}.
For that ceiling it gives the following exact identity.

\begin{proposition}
\label{prop:join}

For every nonempty finite simple graph $G$ and every integer $s\geq0$,

\nopagebreak[4]
\[
\kappa^*(K_s\vee G)=s+\kappa^*(G).                          \tag{J1}
\]

Moreover,

\nopagebreak[4]
\[
\delta(K_s\vee G)=s+\delta(G).                              \tag{J2}
\]

\end{proposition}

\begin{proof}

Choose a minor $F$ of $G$ with $\kappa(F)=\kappa^*(G)$. The graph
$K_s\vee F$ is a minor of $K_s\vee G$, and the join-connectivity formula gives

\nopagebreak[4]
\[
\kappa(K_s\vee F)=s+\kappa(F).
\]

This proves the lower bound in (J1).

For the reverse inequality, take a nonempty minor $H$ of $K_s\vee G$ and a
branch-set model for it. Retain every touching edge between branch sets, thereby
obtaining a supergraph $H^+$ of $H$. Suppose that $r$ branch sets meet the
$K_s$ side. Then $r\leq s$, each of these branch sets is universal in $H^+$,
and the remaining branch sets form a full branch-set quotient $F$ in $G$. If
$F$ is nonempty, then

\nopagebreak[4]
\[
\kappa(H)\leq\kappa(H^+)
=r+\kappa(F)
\leq s+\kappa^*(G).
\]

If $F$ is empty, then $H^+=K_r$ with $r\geq1$, so
$\kappa(H)\leq r-1\leq s+\kappa^*(G)$. Maximizing over $H$ proves (J1).

For (J2), old vertices acquire $s$ new neighbours, while each new clique vertex
has degree $s-1+|V(G)|$. Since $\delta(G)\leq|V(G)|-1$, their minimum is
$s+\delta(G)$.
\end{proof}


Applying Proposition~\ref{prop:join} to (3) gives

\nopagebreak[4]
\[
\delta(K_s\vee M_{12})=s+5,
\qquad
\kappa^*(K_s\vee M_{12})=s+4.                              \tag{J3}
\]

Thus the minimum degree in (J3) is unbounded while the deficit stays equal to
one. This already disproves the existence of a threshold in (4). The remainder
of the paper proves the stronger unbounded-deficit statement.

\section{The graph \texorpdfstring{$M_{12}$}{M12} and its independent blow-ups}


We give a labelled definition so that the construction does not depend on a
drawing. Take two disjoint copies of $C_4\vee K_2$. In the first copy, let the
cycle vertices be $x_0,x_1,x_2,x_3$ in cyclic order and let $a_x,b_x$ be the
two adjacent vertices joined to the entire cycle. In the second copy use
$y_0,y_1,y_2,y_3,a_y,b_y$. Add the six cross edges

\nopagebreak[4]
\[
x_0y_i\quad\text{and}\quad y_0x_i
\qquad (i=1,2,3).                                            \tag{7}
\]

The resulting graph is $M_{12}$. We use the numerical labels

\nopagebreak[4]
\[
\begin{aligned}
0,1,2,3,4,5&=x_0,x_1,x_2,x_3,a_x,b_x,\\
6,7,8,9,10,11&=y_0,y_1,y_2,y_3,a_y,b_y.
\end{aligned}                                                \tag{8}
\]

The graph has 12 vertices and 32 edges. Its vertices $0$ and $6$ have degree
seven and all other vertices have degree five. The set $\{0,6\}$ is a separator;
after its deletion, the two components have vertex sets $\{1,2,3,4,5\}$ and
$\{7,8,9,10,11\}$.

For an integer $t\geq1$, the independent blow-up
$G_t=M_{12}[\overline K_t]$ has vertex set
$V(M_{12})\times\{0,\ldots,t-1\}$. Distinct vertices $(u,i)$ and $(v,j)$ are
adjacent precisely when $uv\in E(M_{12})$. Thus

\nopagebreak[4]
\[
|V(G_t)|=12t,
\qquad
|E(G_t)|=32t^2,
\qquad
\delta(G_t)=5t.                                              \tag{9}
\]

The $t$ vertices above one base vertex will be called a fibre. We write
$I_m=\overline K_m$ for the independent graph on $m$ vertices.

\section{Two elementary connectivity tools}


For a graph $J$, define

\nopagebreak[4]
\[
\beta(J)=\max\bigl(\{0\}\cup\{|X|+|Y|:
X,Y\ne\varnothing,\ X\cap Y=\varnothing,\ E_J(X,Y)=\varnothing\}\bigr).
                                                                    \tag{10}
\]

Thus $\beta(J)=0$ for every complete graph, including the empty graph and the
one-vertex graph. The next formula will only be used when $J$ is not complete.

\begin{lemma}
\label{lem:beta}

If $J$ is a noncomplete graph on $n$ vertices, then

\nopagebreak[4]
\[
\kappa(J)=n-\beta(J).                                       \tag{11}
\]

\end{lemma}

\begin{proof}

If $X$ and $Y$ are nonempty anticomplete sets, deleting all other vertices
disconnects $J$, so $\kappa(J)\leq n-|X|-|Y|$. Conversely, delete a minimum
vertex cut and partition the components that remain into two nonempty unions
$X$ and $Y$. These sets are anticomplete and have total size
$n-\kappa(J)$. Maximizing over $X,Y$ gives (11).
\end{proof}


Recall that the join $J_1\vee J_2$ is obtained from disjoint copies of $J_1$
and $J_2$ by adding every edge between them. In a join of several graphs,
anticomplete shores must lie in a single factor. Consequently, if

\nopagebreak[4]
\[
J=K_u\vee J_1\vee\cdots\vee J_m
\]

is noncomplete, then

\nopagebreak[4]
\[
\kappa(J)=u+\sum_{i=1}^m |V(J_i)|-\max_i\beta(J_i).          \tag{12}
\]

Given pairwise disjoint connected branch sets in a graph $F$, their \textbf{full branch-set quotient} has one vertex per branch set and an edge for every pair
of branch sets that touch in $F$. It is a minor of $F$ obtained without a final
edge-deletion step. The distinction matters: the following profiles are false
for arbitrary minors after unrestricted edge deletion, but hold for the full
quotients used in our upper-bound argument. For connected noncomplete graphs,
the identity (11), in an equivalent complement-biclique form, also appears in
Barrett, Fallat, Hall and Hogben \cite{barrett2013}.

We need three elementary full-quotient profiles.

\begin{lemma}
\label{lem:profiles}

Let $t,B\geq1$ be integers, and let $J_i$ be a full branch-set quotient of the
indicated graph.

\begin{enumerate}
\item If $J_1$ comes from $K_t\sqcup I_t$, then either $J_1$ is complete and
   $|V(J_1)|\leq t$, or $\beta(J_1)=|V(J_1)|$.
\item If $J_2$ comes from $I_{2t}$ and $\beta(J_2)\leq B$, then
   $|V(J_2)|\leq B$.
\item If $J_3$ comes from $K_t\sqcup K_{t,t}$, then one of the following holds:
   $J_3$ is empty; both components contribute vertices and
   $\beta(J_3)=|V(J_3)|$; only the $K_t$ component contributes and $J_3$ is
   complete of order at most $t$; or only the $K_{t,t}$ component contributes
   and

\nopagebreak[4]
\[
   |V(J_3)|\leq\min\{2t,t+B\}                               \tag{13}
\]

   whenever $\beta(J_3)\leq B$.

\end{enumerate}

\end{lemma}

\begin{proof}

A branch set cannot join different components of a graph. Because every
touching edge is retained, a full quotient of $K_t\sqcup I_t$ is a clique
together with isolated vertices. If it is not complete, its vertices can be
partitioned into two nonempty anticomplete shores, proving the first assertion.
The second assertion is immediate because a full quotient of an independent
graph is independent.

For the third assertion, the empty case is immediate, and only the last case
needs proof. We have

\nopagebreak[4]
\[
\mathrm{tw}(K_{t,t})=t.                                    \tag{14}
\]

Taking one bipartition class together with each vertex of the other class
gives a tree decomposition of width $t$. In the opposite direction,
contracting $t-1$ disjoint cross edges and retaining the two unmatched
vertices gives a $K_{t+1}$ minor. Treewidth is minor-monotone, and vertex
connectivity is at most treewidth. Hence every minor of $K_{t,t}$ has
connectivity at most $t$, and every complete minor has order at most $t+1$.

If $J_3$ is noncomplete, Lemma~\ref{lem:beta} gives

\nopagebreak[4]
\[
|V(J_3)|-\beta(J_3)=\kappa(J_3)\leq t.
\]

If $J_3$ is complete, then $|V(J_3)|\leq t+1\leq t+B$.
Together with the trivial bound $|V(J_3)|\leq2t$, this proves (13).

\end{proof}


\section{A uniform lower-bound model}


Fix $t\geq2$ and put

\nopagebreak[4]
\[
a=\left\lfloor\frac t2\right\rfloor,
\qquad
b=\left\lceil\frac t2\right\rceil.                         \tag{15}
\]

Using the labels in (8), form the following $t$ connected branch sets in
$G_t$:

\nopagebreak[4]
\[
\begin{aligned}
A_i&=\{(0,i),(6,i),(7,i),(10,i),(11,i)\}
&& (0\leq i<a),\\
B_j&=\{(1,j),(6,a+j)\}
&& (0\leq j<b).
\end{aligned}                                                \tag{16}
\]

Add as singleton branch sets all unused vertices in the six fibres above
$0,1,2,3,4,5$. There are $t$ branch sets in (16) and $5t$ singletons, hence
$6t$ branch sets in total.

The sets $A_i$ are connected by the edges represented by (7) and the two
copies of $C_4\vee K_2$; each $B_j$ is connected by the edge $1\,6$. The sets
are pairwise disjoint. Moreover, every branch set in (16) is adjacent to every
other branch set in the quotient. Thus the corresponding $t$ quotient
vertices are universal.

Let $Q_t$ be the full branch-set quotient. In the complement of $Q_t$, the
$t$ universal vertices are isolated. The four nontrivial components have
orders

\nopagebreak[4]
\[
t,\qquad t,\qquad 2t-a,\qquad 2t-b.                          \tag{17}
\]

Each component is complete. The two larger components join the remaining
vertices above $0$ to the fibre above $2$, and the remaining vertices above
$1$ to the fibre above $3$; the other two components are the fibres above $4$
and $5$. Hence the largest complete bipartite subgraph in $\overline{Q_t}$ has
total shore size

\nopagebreak[4]
\[
\max\{t,2t-a,2t-b\}=\left\lceil\frac{3t}{2}\right\rceil.
\]

Equivalently, the maximum total size of two anticomplete shores in $Q_t$ is
$\lceil3t/2\rceil$. Lemma~\ref{lem:beta} now gives

\nopagebreak[4]
\[
\kappa(Q_t)
=6t-\left\lceil\frac{3t}{2}\right\rceil
=\left\lfloor\frac{9t}{2}\right\rfloor.                    \tag{18}
\]

Since $Q_t$ is a minor of $G_t$,

\nopagebreak[4]
\[
\kappa^*(G_t)\geq\left\lfloor\frac{9t}{2}\right\rfloor.    \tag{19}
\]

\section{Localization to a torso}


Let $S_t$ be the $2t$ vertices in the fibres above $0$ and $6$. The graph
$G_t-S_t$ has two components of order $5t$. Fix the component above
$\{1,2,3,4,5\}$ and call it $C_t$. Let $T_t$ be obtained from
$G_t[C_t\cup S_t]$ by completing $S_t$ to a clique. Thus $T_t$ has $7t$
vertices.

\begin{lemma}
\label{lem:torso}

If $G_t$ has an $r$-connected minor with $r>2t$, then $T_t$ has an
$r$-connected minor.

\end{lemma}

\begin{proof}

Let $H$ be an $r$-connected minor of $G_t$, represented by pairwise disjoint
connected branch sets. Let $X\subseteq V(H)$ index the branch sets that meet
$S_t$. Since the branch sets are disjoint,

\nopagebreak[4]
\[
|X|\leq|S_t|=2t<r.
\]

The graph $H-X$ is connected and has at least two vertices. Every branch set
indexed by $V(H)\setminus X$ lies in one component of $G_t-S_t$. As there are
no edges between those two components, all these branch sets lie on the same
side; by symmetry, take that side to be $C_t$.

Restrict each branch set meeting $S_t$ to $C_t\cup S_t$. The restricted set is
nonempty. Any excursion through the discarded component enters and leaves at
vertices of $S_t$, so an added clique edge in $T_t[S_t]$ replaces that
excursion. The same clique restores any adjacency between two restricted
branch sets that was realized through the discarded side. Adjacencies to pure
branch sets in $C_t$ remain present. The restricted sets therefore form an
$H$-minor model in $T_t$.
\end{proof}


Since $\lfloor9t/2\rfloor+1>2t$, Lemma~\ref{lem:torso} shows that the upper bound in Theorem
1 follows once we prove

\nopagebreak[4]
\[
\kappa^*(T_t)\leq\left\lfloor\frac{9t}{2}\right\rfloor.     \tag{20}
\]

\section{The three-factor torso}


Partition the seven fibres of $T_t$ into

\nopagebreak[4]
\[
F_1=\{0,2\},
\qquad
F_2=\{1,3\},
\qquad
F_3=\{6,4,5\}.
\]

The definition of $M_{12}$ and the completion of $S_t$ give

\nopagebreak[4]
\[
\begin{aligned}
F_1&\cong K_t\sqcup I_t,\\
F_2&\cong I_{2t},\\
F_3&\cong K_t\sqcup K_{t,t},
\end{aligned}
\]

and every vertex in one factor is adjacent to every vertex in either other
factor. Therefore

\nopagebreak[4]
\[
T_t=(K_t\sqcup I_t)\vee I_{2t}\vee(K_t\sqcup K_{t,t}).      \tag{21}
\]

Consider an arbitrary minor model in $T_t$ and retain every edge of its full
branch-set quotient. Adding these edges cannot decrease vertex connectivity,
so it suffices to bound the connectivity of the full quotient.

Call a branch set mixed if it meets at least two of the three factors in (21),
and pure otherwise. Every mixed branch set becomes a universal quotient
vertex. Indeed, given another branch set, one of the factors met by the mixed
set differs from a factor met by the other set, and a join edge supplies the
required adjacency. The same argument shows that two mixed branch sets are
adjacent.

Let $u$ be the number of mixed branch sets. For $i=1,2,3$, the pure branch sets
contained in $F_i$ form a full branch-set quotient $H_i$ of that factor, and in
particular a minor of it; write $n_i=|V(H_i)|$. The full quotient is

\nopagebreak[4]
\[
H=K_u\vee H_1\vee H_2\vee H_3.                             \tag{22}
\]

Each mixed branch set uses at least two vertices, and each pure branch set uses
at least one. Since $|V(T_t)|=7t$,

\nopagebreak[4]
\[
2u+n_1+n_2+n_3\leq7t.                                     \tag{23}
\]

\section{The general upper bound}


We now prove (20). First suppose that the quotient $H$ in (22) is not complete.
Set

\nopagebreak[4]
\[
B=\max_{1\leq i\leq3}\beta(H_i),
\qquad
s=n_1+n_2+n_3.
\]

Then $B\geq1$, and (12), (22) and (23) give

\nopagebreak[4]
\[
\begin{aligned}
\kappa(H)
&=u+s-B\\
&\leq\frac{7t-s}{2}+s-B
=\frac{7t}{2}+\frac{s}{2}-B.
\end{aligned}
\tag{24}
\]

We distinguish the possible contributions to $H_3$.

\paragraph{Case 1: both components of $F_3$ contribute.}


Lemma~\ref{lem:profiles} gives $B\geq n_3$. Hence

\nopagebreak[4]
\[
\kappa(H)\leq u+n_1+n_2.
\]

Every mixed branch set uses at least one vertex of $F_1\cup F_2$, and each
pure branch set counted by $n_1+n_2$ uses a different vertex there. Since
$|F_1\cup F_2|=4t$,

\nopagebreak[4]
\[
u+n_1+n_2\leq4t
\leq\left\lfloor\frac{9t}{2}\right\rfloor.                 \tag{25}
\]

\paragraph{Case 2: only the $K_t$ component of $F_3$ contributes, or $H_3$ is empty.}


Here $n_3\leq t$. If $H_1$ is noncomplete, Lemma~\ref{lem:profiles} gives
$n_1,n_2\leq B$, and therefore $s\leq t+2B$. Equation (24) yields

\nopagebreak[4]
\[
\kappa(H)\leq4t.                                           \tag{26}
\]

If $H_1$ is complete, then $n_1\leq t$ and $n_2\leq B$, so
$s\leq2t+B$. Hence

\nopagebreak[4]
\[
\kappa(H)\leq\frac{9t}{2}-\frac B2
\leq\left\lfloor\frac{9t}{2}\right\rfloor.                \tag{27}
\]

The final inequality uses $B\geq1$ and the integrality of $\kappa(H)$.

\paragraph{Case 3: only the $K_{t,t}$ component of $F_3$ contributes.}


Suppose first that $H_1$ is noncomplete. Lemma~\ref{lem:profiles} gives

\nopagebreak[4]
\[
s\leq2B+\min\{2t,t+B\}.                                    \tag{28}
\]

If $B\leq t$, then $s\leq t+3B$, and (24) gives

\nopagebreak[4]
\[
\kappa(H)\leq4t+\frac B2\leq\frac{9t}{2}.
\]

If $B\geq t$, then $s\leq2t+2B$, and (24) again gives
$\kappa(H)\leq9t/2$.

It remains to suppose that $H_1$ is complete. Now

\nopagebreak[4]
\[
s\leq t+B+\min\{2t,t+B\}.                                  \tag{29}
\]

For $B\leq t$, this gives $s\leq2t+2B$ and
$\kappa(H)\leq9t/2$. For $B\geq t$, it gives $s\leq3t+B$ and

\nopagebreak[4]
\[
\kappa(H)\leq5t-\frac B2\leq\frac{9t}{2}.
\]

In every subcase, integrality yields

\nopagebreak[4]
\[
\kappa(H)\leq\left\lfloor\frac{9t}{2}\right\rfloor.       \tag{30}
\]

The three cases are mutually exclusive and exhaustive for a noncomplete
quotient.

Finally suppose that $H$ is complete. Every nonempty $H_i$ must be complete.
Lemma~\ref{lem:profiles} gives

\nopagebreak[4]
\[
n_1\leq t,
\qquad
n_2\leq1,
\qquad
n_3\leq t+1,
\]

and thus $s\leq2t+2$. By (23),

\nopagebreak[4]
\[
|V(H)|=u+s
\leq\frac{7t+s}{2}
\leq\frac{9t+2}{2}.
\]

Since $\kappa(H)=|V(H)|-1$, the same parity check gives

\nopagebreak[4]
\[
\kappa(H)\leq\left\lfloor\frac{9t}{2}\right\rfloor.       \tag{31}
\]

Equations (30) and (31) prove (20). Lemma~\ref{lem:torso} transfers this upper bound to
$G_t$, and (19) supplies equality. This completes the proof of Theorem~\ref{thm:main}.

\section{Consequences}


\begin{corollary}
\label{cor:unbounded}

There is no constant $C$ such that

\nopagebreak[4]
\[
\kappa^*(G)\geq\delta(G)-C
\]

for every nonempty finite simple graph $G$.

\end{corollary}

\begin{proof}

For $G_t$, the deficit is $\lceil t/2\rceil$ by (6), and hence is unbounded.

\end{proof}


\begin{corollary}
\label{cor:threshold}

There is no constant $c$ such that every nonempty graph $G$ with
$\delta(G)\geq c$
satisfies $\kappa^*(G)\geq\delta(G)$.

\end{corollary}

\begin{proof}

This already follows from (J3). It also follows from the stronger family
$G_t$: the minimum degrees $\delta(G_t)=5t$ are unbounded, while (6) is
positive for every $t\geq2$.
\end{proof}


For $t\geq3$, let $k=5t-1$. Then $G_t$ has minimum degree $k+1$ but

\nopagebreak[4]
\[
\kappa^*(G_t)=\left\lfloor\frac{9t}{2}\right\rfloor<5t-1=k.
\]

Thus the family gives infinitely many counterexamples to Barát's conjecture.
The first member covered by this strict inequality is the 36-vertex graph
$G_3$, for which $\delta(G_3)=15$ and $\kappa^*(G_3)=13$.

The formula also gives $\kappa^*(G_2)=9$ and $\kappa^*(G_4)=18$. The previously
known value $\kappa^*(M_{12})=4$ is consistent with the same rounded expression
at $t=1$, although Theorem~\ref{thm:main} is stated only for $t\geq2$.

\section{Reproducibility and evidence boundary}


The proof of Theorem~\ref{thm:main} is independent of exhaustive computation. The
accompanying verification program reconstructs the branch sets in (16), checks
that they are nonempty, disjoint and connected in the original graph, builds
their full quotient, and computes its vertex connectivity. From the repository
root, run

\begin{quote}\small

\noindent{\small\ttfamily .\textbackslash{}\allowbreak{}.venv\textbackslash{}\allowbreak{}Scripts\textbackslash{}\allowbreak{}python.exe\ \allowbreak{}research\textbackslash{}\allowbreak{}M12-Minor-Connectivity-Gap\textbackslash{}\allowbreak{}verification\textbackslash{}\allowbreak{}verify\_general\_lower\_model.py\par}
\end{quote}


The frozen output covers $t=2,\ldots,8$ and reports the connectivities

\nopagebreak[4]
\[
9,13,18,22,27,31,36.
\]

This finite calculation checks the stated construction and its implementation;
it is not used to infer the universal quantifier in Theorem~\ref{thm:main}. The upper bound
is entirely theoretical.

\section{Prior work and novelty boundary}


Barát's thesis records the conjecture leading to (2) \cite{barat}. The graph $M_{12}$
originates in Mader's 1968 work \cite{mader1968}; Barioli et al. later introduced the relevant
minor-monotone ceiling notation and established (3) for this graph \cite{barioli2013}. Barrett,
Fallat, Hall and Hogben developed nearby connectivity tools \cite{barrett2013},
including an equivalent form of (11) and the ordinary join-connectivity
formula. Fijav\v{z} and Wood studied broader relations between graph minors and
minimum degree \cite{fijavz2010}. Barioli, Fallat, Gupta and Li subsequently used $M_{12}$ \cite{barioli2026} to
show that $\kappa^*(G)\geq\delta(G)$ does not hold for every graph and asked
whether it holds above a universal minimum-degree threshold; Proposition~\ref{prop:join}
gives a negative answer, while Theorem~\ref{thm:main} supplies the new independent-blow-up
formula, makes the deficit unbounded, and yields infinitely many counterexamples
to Barát's conjecture.

A targeted search did not locate the exact formula (5) or this independent
blow-up family in prior work. This is a candidate novelty statement rather
than proof of global priority. No claim of first discovery is made, and the
historical assessment remains subject to specialist review and searches of
sources not indexed in the databases consulted.

\section*{AI-assisted research and writing disclosure}


Carptopus is the sole author and is responsible for the mathematical claims,
source selection, revisions and public version of this manuscript. OpenAI
Codex was used as an AI-assisted tool for literature discovery, proof
exploration, exact verification, adversarial auditing and manuscript drafting.

\begin{thebibliography}{9}
\footnotesize
\raggedright

\bibitem{barat}
J.~Bar\'at,
\emph{Width-type graph parameters},
PhD thesis, University of Szeged, submitted 2000 and defended 2001,
Conjecture 6.2.
\url{https://doktori.bibl.u-szeged.hu/id/eprint/4651/}.

\bibitem{barioli2013}
F.~Barioli, W.~Barrett, S.~M.~Fallat, H.~T.~Hall, L.~Hogben, B.~Shader,
P.~van den Driessche, and H.~van der Holst,
\emph{Parameters related to tree-width, zero forcing, and maximum nullity of a graph},
Journal of Graph Theory \textbf{72} (2013), no.~2, 146--177.
\href{https://doi.org/10.1002/jgt.21637}{doi:10.1002/jgt.21637}.

\bibitem{barioli2026}
F.~Barioli, S.~M.~Fallat, H.~Gupta, and Z.~Li,
\emph{The weak version of the graph complement conjecture and partial results for the delta conjecture},
Discrete Mathematics \textbf{349} (2026), no.~3, 114861.
\href{https://doi.org/10.1016/j.disc.2025.114861}{doi:10.1016/j.disc.2025.114861}.

\bibitem{barrett2013}
W.~Barrett, S.~M.~Fallat, H.~T.~Hall, and L.~Hogben,
\emph{Note on Nordhaus--Gaddum problems for Colin de Verdi\`ere type parameters},
Electronic Journal of Combinatorics \textbf{20} (2013), no.~3, P56.
\url{https://www.combinatorics.org/ojs/index.php/eljc/article/view/v20i3p56}.

\bibitem{mader1968}
W.~Mader,
\emph{Homomorphies\"atze f\"ur Graphen},
Mathematische Annalen \textbf{178} (1968), 154--168.
\href{https://doi.org/10.1007/BF01350657}{doi:10.1007/BF01350657}.

\bibitem{fijavz2010}
G.~Fijav\v{z} and D.~R.~Wood,
\emph{Graph minors and minimum degree},
Electronic Journal of Combinatorics \textbf{17} (2010), R151.
\href{https://doi.org/10.37236/423}{doi:10.37236/423}.

\end{thebibliography}

\end{document}
